\section{Unary contexts and the meaning of repetition}\label{app:contexts}
This appendix classifies the one-variable contexts of the reduct, explains which observations retain ambiguity, and identifies a restriction under which Nash and subgame-perfect equilibria have the same possible winners. It separates bottom from omitted contributions and relates collective decision rules through their outcome interpretations. It then separates three readings of a repeated component: independent copies, uniform equilibrium selection, and a shared hidden decision.

\subsection{Unary contexts without internal switches}

Assume $K\geq3$ and the reduct $\mathbf A$, without internal switches.
A unary term function uses one variable, any number of
its occurrences, owned choices, and leaf constants. By
\cref{game:thm:union}, arbitrary nonempty ground constants are available.
Write $\Reg(C,E)$ for the following function, with parameters
\[
 C\subseteq I,\qquad E_i\subseteq C\setminus\{i\}\quad(i\in I):
\]
\begin{equation}\label{game:eq:normal}
 \Reg(C,E)(S)=
 \begin{cases}
 \{i\}\cup E_i,&S=\{i\},\\
 S\cup C,&|S|\geq2.
 \end{cases}
\end{equation}

\begin{proposition}[Unary normal forms]\label{game:thm:unary}
The unary term functions are exactly the nonempty constant functions and
the functions~\eqref{game:eq:normal}. The regular parameters $C,(E_i)$ are
unique. The only overlap with the constant functions is the constant
function $I$.
\end{proposition}
\begin{proof}
\emph{Upper bound.} The variable is regular with $C=\varnothing$ and
empty rows, and each constant is constant. Suppose two child functions
are regular with extra sets $C,D$. On a multiwinner input $S$, each child
contains $S$ and cannot equal an owner singleton. Their owned combination
is their union, $S\cup C\cup D$. On singleton $\{i\}$, both contain $i$;
their owned combination still contains $i$ and is contained in
$\{i\}\cup C\cup D$. This is regular.

If one child is constant $D$ and the other regular with extra set $C$,
then $D=\{r\}$ makes the result constant. Otherwise the regular child's
multiwinner values cannot collapse and the result is $S\cup C\cup D$.
On an atomic input, the operation preserves $i$ and stays inside
$\{i\}\cup C\cup D$. Two constants give a constant. This proves closure
and the upper classification by induction.

\emph{Synthesis.} For every $k\in C$, prescribe the allowed singleton
input roles
\[
 A_k=\{i\ne k:k\in E_i\}.
\]
Start with $g_k(x)=J(x,\at k)$. For each $i\ne k$ outside $A_k$, replace
$g_k(x)$ by $g_k(x)\ow i x$. At singleton input $i$, this suppresses the
extra winner $k$ by returning $\{i\}$. At another singleton, it preserves
the current value, which contains that other role. On a multiwinner
input $S$, both operands contain $S$, so the result remains $S\cup\{k\}$.
The order of these suppressions is irrelevant. Join $x$ with all the
$g_k(x)$ using $J$. This realizes precisely~\eqref{game:eq:normal}; for empty
$C$, use $x$. Its terms have only
singleton leaves and owned choices.

\emph{Uniqueness.} For any $k$, choose a two-element $S$ omitting $k$.
Then $k\in C$ exactly when $k\in\Reg(C,E)(S)$, so $C$ is uniquely
determined. Each $E_i$ is the singleton-input result minus $i$. A regular
function is extensive. If it is constant, its value must contain every
singleton and hence equal $I$. That constant is regular with $C=I$ and
$E_i=I\setminus\{i\}$.
\end{proof}

\begin{corollary}[Count and a small test set]
There are exactly
\begin{equation}\label{game:eq:count}
  (1+2^{K-1})^K+2^K-2
\end{equation}
unary term functions. Equality of two unary functions is determined by
their values on singleton and two-element input sets.
\end{corollary}
\begin{proof}
For fixed $C$ there are $2^{|C|(K-1)}$ row families. Summing over $C$
gives $(1+2^{K-1})^K$. Add the $2^K-1$ constants and subtract the single
overlap. Singleton/pair tests determine all regular parameters. They
also distinguish constants: a constant agreeing with a regular function
on every singleton must contain every role, hence be $I$; pair tests
then force $C=I$ and all rows to be full.
\end{proof}

Thus there are $131$ unary functions for $K=3$, $6575$ for $K=4$, and
$1419887$ for $K=5$. The first number is independently recovered by
exhaustive algebraic closure in \texttt{check\_choice.py}. The script constructs
every parameterized function by actual owned-choice operations for
$K=3,4$ and checks separation by singleton/pair tests. It does not
enumerate all terms for $K=4,5$; the general classification supplies the
upper bound.

The parameters are also closed under context composition. If the outer
function has parameters $(C,E)$ and the inner $(D,F)$, their composite is
regular with extra set $C\cup D$ and rows
\[
 H_i=\begin{cases}
 E_i,&F_i=\varnothing,\\
 (F_i\cup C)\setminus\{i\},&F_i\ne\varnothing.
 \end{cases}
\]
If either function is constant, evaluate the appropriate constant case.
Indeed, a multiwinner input stays multiwinner through both functions. At
input $\{i\}$, the inner result is still a singleton exactly when
$F_i=\varnothing$; otherwise the outer function adds $C$.
This supplies an explicit finite monoid of unary contexts.

Atomic inputs alone do not determine a context. For $I=\{a,b,c\}$, the term
\begin{equation}\label{game:eq:D}
 D(x)=x\ow c\bigl((x\ow a\{b\})\ow b\{a\}\bigr)
\end{equation}
fixes every singleton but sends each multiwinner set $S$ to
$S\cup\{a,b\}$. It has only two occurrences of $x$. Thus a calculus that
checks equations only on the winner constants misses real distinctions.

\subsection{Which observations retain the information?}

\begin{proposition}[Unique-winner tests collapse ambiguity]\label{game:prop:obs}
Every unary context function is either constant or extensive. If its
observation is only whether its result equals a designated singleton,
all sets with at least two members are contextually indistinguishable.
The family of membership observations, with the observing role allowed
to vary, distinguishes distinct sets already through the identity context.
\end{proposition}
\begin{proof}
The first statement follows either directly by induction (a constant
owner-singleton absorbs; otherwise a nonconstant child is preserved), or
from \cref{game:thm:unary}. A constant context has the same result on all
inputs. An extensive context cannot turn a multiwinner input into a
singleton. Membership tests distinguish sets at an element of their
symmetric difference. One fixed role need not distinguish a pair of inputs: for example, $\{k\}$
and $\{k,j\}$ both contain $k$.
The exact fixed-role contextual quotient is \cref{view:thm:view}.
\end{proof}

Restricting observations to unique
winners can therefore discard the ambiguity retained by possible-win
membership tests.

For comparison, let $q(S)=\{i:S=\{i\}\}$. For nonempty operands,
\[
 q(S\ow r T)_r=q(S)_r\lor q(T)_r,\qquad
 q(S\ow r T)_i=q(S)_i\land q(T)_i\quad(i\ne r).
\]
These are precisely the coordinate operations in the guarantee algebra
of Olde Loohuis and Venema, Definition 3.20
\cite{OldeLoohuisVenema2010}, restricted to its empty and singleton
values. Induction also shows that, for every public game $G$, $\den{G}=\{i\}$
exactly when $i$ can force a win. This comparison explains the
collapse; forcing is not how $\den{G}$ is defined.

The full algebra validates different equations. For example, with
$r=0,s=1$, $A=\{0,1\}$, $B=\{1\}$, and $C=\{2\}$,
\[
 A\ow0(B\ow1 C)=\{0,1\},\qquad
 (A\ow0 B)\ow1(A\ow0 C)=\{0,1,2\}.
\]
Thus the mutual distributivity axiom of their multiplayer algebras
fails here, even though every $\star_r$ distributes over the neutral
union operation. Reusing their complete equational calculus would
therefore change the semantics.

\subsection{Choices paired with distinct winning leaves}\label{app:contexts-unrepeated}

A finite binary tree with $n$ choice nodes has $n+1$ leaves, so its
individual choices and leaves cannot be in bijection. More generally,
a finite tree whose choice nodes all have at least two children has
more leaves than choices. The closest binary case assigns every choice
a different winning leaf and leaves one leaf with no choices of its
own. In the standard presentation, each controller then owns one
choice and each outcome occurs at one leaf. This case admits a simple
structural calculation, and a larger class has the same simplification.

Say that ownership is \emph{unrepeated on plays} if no root-to-leaf
path contains two choices owned by the same role. A role may still own
choices in separate branches. For a subtree $H$, write $L(H)$ for
the set of all its leaf winners, before any equilibrium selection.

\begin{proposition}[Unrepeated ownership]\label{context:prop:unrepeated}
For a finite public game with ownership unrepeated on plays, define
\[
 P(\at i)=\{i\},\qquad
 P(G\ch r H)=
 \begin{cases}
 \{r\},&L(G)=\{r\}\text{ or }L(H)=\{r\},\\
 P(G)\cup P(H),&\text{otherwise}.
 \end{cases}
\]
Then $P(G)$ is both the set of possible SPE winners and the set of
possible Nash winners, for the paper's one-winner payoffs.
\end{proposition}
\begin{proof}
If a subtree contains no choices owned by $r$, its SPE value equals
$\{r\}$ exactly when all its leaves are won by $r$. Indeed, a choice
owned by $s\ne r$ either returns $\{s\}$ or takes union. Returning
$\{r\}$ therefore requires both children to return $\{r\}$, so the
claim follows inductively; the converse is immediate.

At a node owned by $r$, neither child contains an $r$-choice. The
owner-singleton test in backward induction is consequently exactly
$L(G)=\{r\}$ or $L(H)=\{r\}$. Thus the displayed recursion computes
the possible SPE winners.

Every SPE is Nash. Conversely, restrict a Nash profile to its reached
child; this is a Nash profile of that child, since any profitable
deviation there extends to the whole game while retaining the prefix.
By induction, its winner belongs to that child's $P$ value. If neither
child is monochromatic for the root owner, it therefore belongs to
the parent's union. If a child is monochromatic for $r$, the root
owner can win by selecting it regardless of all subsequent actions,
so every Nash profile is won by $r$. This proves the other inclusion.
\end{proof}

\begin{corollary}[Injective matching is local pruning]\label{context:cor:injective-matching}
Suppose each choice is paired with a different winning leaf, and
different leaves have different winners. The possible Nash and SPE
winners are the leaves remaining after the following pruning: whenever
a choice has its own winning leaf as an immediate child, replace that
choice subtree by the winning leaf; retain both children at every
other choice. In particular, a game with an internal root has a unique
equilibrium winner exactly when its root is paired with an immediate
leaf child.
\end{corollary}
\begin{proof}
Distinct owners give unrepeated ownership. With distinct leaf winners,
a child subtree is monochromatic for the parent owner exactly when it
is that owner's leaf. Thus the recursion for $P$ is precisely this
pruning. If the root is not pruned, its two child values are nonempty
and disjoint, so its value is ambiguous.
\end{proof}

The equality concerns possible winners, not complete profiles. For
example, in $\at a\ch a(\at b\ch b\at c)$, choosing $a$ at the
root and prescribing $c$ on the right is Nash but not subgame perfect.
The unique winner is $a$ under either solution concept. A separation
of their possible-winner observations therefore requires a repeated
owner on some play in the present payoff model.

\subsection{Bottom, top, and omitted subtrees}\label{app:contexts-bottom}

An empty contribution to a union, an identity for owned choice, and a
terminal at which everyone loses are three different constructions.
The first already occurs in the status forms of
\cref{view:thm:normal}: either field of $(C,\Lambda)$ may be empty.
Once the singleton tests have been fixed, an absent contribution can
therefore be omitted. The following obstruction concerns putting such
a unit into the game algebra itself.

\begin{proposition}[A common unit removes ownership]\label{context:prop:common-unit}
Let $B$ satisfy all identities of the constant-free core $\mathbf C_K$,
and let $+$ denote its derived union $J$. If $e\star_r x=x$ for every
owner $r$ and every $x\in B$, then $x\star_r y=x+y$ for every $r$.
Consequently no extension containing $\mathbf C_K$ as a subalgebra can
adjoin a common unit while preserving all its identities.
\end{proposition}
\begin{proof}
The core validates the semilattice laws for $+$, commutativity of each
$\star_r$, distribution over $+$, and
$(x\star_r y)+x+y=x+y$. These follow directly from the owner-singleton
rule. A common unit gives $e+y=J(e,y)=y$, hence
\[
 x\star_r y=x\star_r(e+y)=x+(x\star_r y).
\]
Thus $x\leq x\star_r y$ in the join order. Commutativity gives the
same inequality for $y$, while the displayed coverage identity gives
$x\star_r y\leq x+y$. All choices are therefore join. In the original
algebra, $\{a\}\star_r\{r\}=\{r\}\ne\{a,r\}$ for $a\ne r$.
\end{proof}

Keeping the formula~\eqref{game:eq:own} on all subsets does give
$\varnothing\star_r S=S$, and each choice remains a semilattice.
But it breaks distribution: for $a\ne r$,
\[
 \{a\}\star_r(\varnothing\cup\{r\})=\{r\},\qquad
 (\{a\}\star_r\varnothing)\cup(\{a\}\star_r\{r\})=\{a,r\}.
\]
Here $J$ still denotes actual union: with an empty operand all three
choices in $J$ return the other operand, and nonempty operands use
\cref{game:thm:union}. Even a unit used only in proofs would therefore
prevent unrestricted use of the original identities.

An annihilator has a different outcome.
Adjoin $0$ with $0\star_r x=0$ for every $r$. This preserves exactly
the identities in the original constant-free core signature: by
\cref{gen:cor:trace}, both sides of a valid identity have the same
variables, so assigning $0$ to any occurring variable makes both sides
$0$; otherwise the original identity applies. The original algebra is
a subalgebra, so no new identity in that signature becomes valid.
This is the zero extension used in \cref{query:prop:relations}.
Its actual set union has $\varnothing$ as unit, whereas its derived
$J$ has $\varnothing$ as annihilator. Thus the zero extension does not
preserve the additional equation $J(x,y)=x\cup y$ on empty inputs.
It helps the relational representation of one owner, but cannot erase
a branch of an owned-choice term.

Nor is the existing value $I$ a common annihilator:
$I\star_r\{r\}=\{r\}$. It is the greatest value under inclusion,
recording that every role is a possible winner.

\paragraph{A unit after restricting the observer and the owners.}
There is a smaller setting in which a bottom does work. Fix an observer
$k$ and a set $O\subseteq I$ of allowed choice owners. Contexts may
use arbitrary nonempty constants and only the $\star_r$ with $r\in O$;
they have no internal role switches. Write $S\equiv_{k,O}T$ when these
contexts agree on whether $k$ belongs to their output.

\begin{proposition}[Restricted player views]\label{context:prop:restricted-view}
The classes of $\equiv_{k,O}$ are exactly
\[
 \pv:\ k\in S,\qquad
 a_i:\ S=\{i\}\quad(i\in O\setminus\{k\}),\qquad
 \nv:\ \text{all remaining values}.
\]
There are $|O\setminus\{k\}|+2$ classes. If $O$ is nonempty, this
quotient has a common unit for its allowed choices exactly when
$|O|=1$; that unit is $\nv$.
\end{proposition}
\begin{proof}
For owner $k$, output membership of $k$ is disjunction of input
membership. For owner $r\ne k$, it is that disjunction provided
neither input is $\{r\}$. An output equals an active opponent atom
$\{i\}$ exactly when either the owner is $i$ and an input is
$\{i\}$, or both inputs are $\{i\}$. Thus the displayed information
is preserved by every allowed operation. Membership separates $\pv$;
the context $x\star_i\{k\}$ separates each $a_i$ from every other
negative class. All classes are inhabited: any two-element set omitting
$k$ represents $\nv$.

A unit must fix $\nv$. The class $\pv$ fails this under every owner,
and $a_i$ fails it under owner $i$, so only $\nv$ is possible. With
one owner, combining $\nv$ with any class returns that class. With
two or more owners, choose $i\in O\setminus\{k\}$ and
$j\in O\setminus\{i\}$. Then
$\nv\star_j a_i=\nv\ne a_i$, excluding a common unit.
\end{proof}

This gives an exact meaning to an omitted subtree for one observer
under one surrounding owner. With more owners, its negative summary
must still distinguish their singleton traps. Having no choices owned
by $k$ inside the subtree is insufficient: the leaves $\at b$ and
$\at c$ both have that property and both lose for $k$, but
$x\star_b\{k\}$ distinguishes them, for distinct $b,c,k$.

\paragraph{Playable losses and simultaneous wins.}
If terminals may declare any winning set $W\subseteq I$, including
$\varnothing$ and $I$, each player still receives payoff one precisely
when it belongs to $W$. An exact value is the nonempty family
$\mathcal F\subseteq\mathcal P(I)$ of terminal winning sets reached
in SPE. Put $\mathcal F_r=\{W\in\mathcal F:r\in W\}$.

\begin{proposition}[Winning-set families]\label{context:prop:outcome-families}
For finite public games with arbitrary terminal winning sets, the SPE
families compose by
\[
 \mathcal F\mathbin{\odot_r}\mathcal H=
 \begin{cases}
 \mathcal F_r\cup\mathcal H_r,
   &\mathcal F_r=\mathcal F\text{ or }\mathcal H_r=\mathcal H,\\
 \mathcal F\cup\mathcal H,&\text{otherwise}.
 \end{cases}
\]
Their unions, which record only possible winning players, do not
determine the union after a choice. This already fails when the only
new terminal is one where everyone loses.
\end{proposition}
\begin{proof}
Fix an SPE in each child. The owner selects an owner-winning outcome
if either prescribed outcome wins for it, and otherwise either branch.
If a child always wins for the owner, only owner-winning outcomes can
be selected, and each such outcome is realized by selecting its child.
Otherwise both children admit owner-losing SPEs. Prescribing one off
the path realizes every outcome in the other child. The whole-player
deviation argument of \cref{game:thm:adequacy} applies unchanged.

For the obstruction, let $d$ be an everyone-loses terminal and take
$G=d\ch b\at a$ and $H=\at a$, with $a\ne b$. Their families are
$\{\varnothing,\{a\}\}$ and $\{\{a\}\}$, whose unions both equal
$\{a\}$. In the context $x\ch a\at b$, the resulting unions are
$\{a,b\}$ and $\{a\}$, respectively. For $G$, prescribing $d$ in
the left child lets $a$ rationally choose the right child. For $H$,
$a$ must choose its certain win.
\end{proof}

With only the extra terminal $d$, these families can be encoded as
nonempty subsets of $I\cup\{d\}$, retaining $d$ as a possible outcome
but never as an owner. This preserves the owned-choice formula, with
$\{d\}$ an ordinary outcome atom rather than a unit. With simultaneous
wins, a family distinguishes $\{I\}$ from
$\{\{i\}:i\in I\}$, although their unions agree. Bottom and top as
playable outcomes therefore require more information than the present
winner sets retain.

In the guarantee algebra of Olde Loohuis and Venema, the same symbols
have a different use. Their Definition~3.20
\cite{OldeLoohuisVenema2010} gives
\[
 \bot\vee_r S=S\cap\{r\},\qquad
 \top\vee_r S=S\cup\{r\},
\]
so bottom isolates one player's guarantee. The unit for $\vee_r$ is
$I\setminus\{r\}$, and its annihilator is $\{r\}$. These are useful
observer operations, but their laws concern guarantees rather than
possible SPE winners.

\subsection{Decision rules, collective choice, and merging roles}\label{app:contexts-collective}

Separate the finite tree, its controllers, and its payoff interpretation.
Let $C$ label the controllers of internal nodes and let $\Omega$ label
terminal outcomes. Payoffs are functions $u_c:\Omega\to\mathbb R$.
The original game is the special case $C=\Omega=I$ with
$u_r(i)=[i=r]$. More generally, a map $w:C\to\Omega$ assigns each
controller a preferred outcome, with $u_c(i)=[i=w(c)]$. Several
controllers may prefer the same outcome; some outcomes may have no
controller preferring them. The tree and the interpretation can
therefore vary separately. This is the distinction between outcome
functions and player selection functions used by Hedges et al.\
\cite[Definitions~1, 3 and 7]{HedgesEtAl2017Selection}.
Distinct controller tags can carry the same payoff function or local
rule; an aggregation can still count their votes separately.

For binary nodes, the local decision rule can be given by a nonempty
set $m_c(x,y)\subseteq\{x,y\}$ of admissible outcomes for each fixed
pair $x,y\in\Omega$. For a maximizing controller, this set consists
of the outcomes attaining $\max(u_c(x),u_c(y))$. Define its lift by
\begin{equation}\label{context:eq:rule-lift}
 S\mathbin{\odot_c}T=\bigcup_{x\in S,\,y\in T}m_c(x,y).
\end{equation}
For utility maximization this is exactly the possible SPE outcome
recurrence: prescribe one child SPE with outcome $x$ and one with
outcome $y$, and select an optimal branch. Other rules specify local
admissibility, with no claim that they are utility-maximizing SPE.

\begin{proposition}[Atomic rules and their transport]\label{context:prop:rule-lift}
Every lift~\eqref{context:eq:rule-lift} is idempotent and distributes
over union in each argument. It is commutative exactly when $m_c$ is
symmetric, and associative exactly when, for every $x,y,z$,
\begin{equation}\label{context:eq:atomic-assoc}
 \bigcup_{v\in m_c(x,y)}m_c(v,z)
 =\bigcup_{v\in m_c(y,z)}m_c(x,v).
\end{equation}
Conversely, every binary operation on nonempty subsets of finite
$\Omega$ that distributes over union and never introduces an outcome
has this form, uniquely.

Let $f:\Omega\to\Omega'$ and $\gamma:C\to C'$. The direct-image
map $S\mapsto f[S]$ transports every operation $\odot_c$ to
$\odot'_{\gamma(c)}$ exactly when
\begin{equation}\label{context:eq:atomic-transport}
 f[m_c(x,y)]=m'_{\gamma(c)}(f(x),f(y))
 \qquad(c\in C,\ x,y\in\Omega).
\end{equation}
\end{proposition}
\begin{proof}
Every outcome in a diagonal product comes from its input, and the
pair $(x,x)$ retains each input outcome. This gives idempotency.
Distribution and the sufficiency of the symmetry and associativity
conditions follow by expanding unions. Necessity follows on singleton
inputs. Conversely, finite distribution expands any proposed operation
as the union of its values on singleton pairs, which uniquely recover
$m_c$. Finally, direct image commutes with unions, so
\eqref{context:eq:atomic-transport} is sufficient for transport of the
lifts, and singleton inputs show it is necessary.
\end{proof}

Maximizing a scalar utility satisfies~\eqref{context:eq:atomic-assoc}:
both sides retain precisely the utility-maximal outcomes among
$x,y,z$. The corresponding lifts are therefore semilattices.
If $u_c=u'_{\gamma(c)}\circ f$, their maximizing rules satisfy
\eqref{context:eq:atomic-transport}. For bijections of controllers and
outcomes this is ordinary relabelling. The original role permutation
uses the same bijection on both sets, preserving the diagonal payoff
table $[i=r]$. Thus role switching transports both the controller and
the outcome interpretation. A coalition can instead aggregate payoff
rows: a scalar rule $h_B$ for $B\subseteq C$ gives the collective utility
\[
 u_B(o)=h_B\bigl((u_c(o))_{c\in B}\bigr).
\]
More general aggregation can directly specify $m_B(x,y)$. Neither
construction requires an invertible relabelling.

\paragraph{A shared binary objective.}
Take $\Omega=I$ and give each allowed coalition $A\subseteq I$ one
controller with payoff $u_A(i)=[i\in A]$. Its nodes have a shared
objective: some member wins. Continuation games retain their own
controllers, and several controllers may value the same label at one.

\begin{proposition}[Choice with a shared objective]\label{context:prop:collective}
For finite public games with these controllers, possible SPE winners
compose by
\begin{equation}\label{context:eq:collective}
 S\star_A T=
 \begin{cases}
 (S\cup T)\cap A,&S\subseteq A\text{ or }T\subseteq A,\\
 S\cup T,&\text{otherwise}.
 \end{cases}
\end{equation}
Every $\star_A$ is associative, commutative and idempotent, and
distributes over union. Moreover,
$\star_{\{r\}}=\star_r$, while $\star_\varnothing$ and
$\star_I$ are both union. If $2\leq|A|<K$, this operation is not
definable by a term of the original full signature
$\mathbf A_K^{\rm sym}$.
\end{proposition}
\begin{proof}
If $S\subseteq A$ or $T\subseteq A$, every pair of child outcomes
offers an $A$-winner, so the maximizing rule keeps precisely the
$A$-winners. Otherwise an outcome outside $A$ in either child makes
every outcome of the other child admissible. This gives the formula
from~\eqref{context:eq:rule-lift}. Its laws follow from
\cref{context:prop:rule-lift} and scalar maximization; the singleton
and endpoint formulas are immediate.

For nondefinability, choose $a\in A$ and $c\notin A$. The operation
is commutative and depends on both inputs. If it had a term in the
original full signature, idempotent elimination
(\cref{gen:thm:core}) would give a core term using both variables.
On ambiguous inputs that term must be their union by
\cref{gen:cor:trace}. But the ambiguous inputs $S=A$ and
$T=\{a,c\}$ give $S\star_A T=A\ne A\cup\{c\}$.
\end{proof}

At a fixed pair of terminal labels, the same rule selects the
Pareto-undominated alternatives for the members' individual payoffs:
a member's win dominates an outsider's win, while wins by different
members are incomparable. Requiring every member to improve strictly
would instead make every alternative admissible when $|A|\geq2$,
since a single terminal label cannot win for two members. The shared
objective therefore specifies a substantive agreement rule.

\paragraph{Majority depends on what is aggregated.}
For a nonempty group $B$ with binary payoffs, a quota rule gives
\[
 u_{B,q}(o)=\left[\sum_{c\in B}u_c(o)\geq q\right].
\]
Quota one is the shared objective, quota $|B|$ is conjunction, and
$q=\lfloor|B|/2\rfloor+1$ is strict majority approval of an outcome.
With the original diagonal payoffs, every quota $q\geq2$ is constantly
zero and its maximizing lift is union. Under a preferred-outcome map
$w$, the sum instead counts $|B\cap w^{-1}(o)|$. For example, two
controllers preferring $a$ and one preferring $b$ give $a$ majority
approval. A general binary payoff table permits several outcomes to
have majority approval.

Comparing alternatives by majority of strict preferences, with
indifferent voters abstaining, uses a different aggregation.
One can also vote on the admissible sets proposed by the tagged
rules. For $1\leq q\leq|B|$, put
\[
 Q_q((X_c)_{c\in B})=\{o:\#\{c\in B:o\in X_c\}\geq q\}.
\]
Quotas above one may leave no admissible outcome.

\begin{lemma}[Aggregation before and after lifting]\label{context:lem:aggregation-lift}
For any finite family of atomic rules and any quota $q$,
\[
 \bigcup_{x\in S,\,y\in T}Q_q((m_c(x,y))_{c\in B})
 \ \subseteq\ Q_q((S\odot_c T)_{c\in B}).
\]
At quota one this is equality; at larger quotas it can be strict.
\end{lemma}
\begin{proof}
A fixed pair witnessing $q$ approvals also witnesses each approval
after lifting. At quota one, the unions over pairs and controllers
commute. For strictness at quota two, take the original rules of
$a,b$, with $S=\{a,b\}$ and $T=\{c\}$. On the left, the pair
$(a,c)$ retains only $a$ in common and $(b,c)$ retains only $b$,
giving $\{a,b\}$. On the right, each delegate admits
$\{a,b,c\}$ using a different witness pair.
\end{proof}

\begin{proposition}[Majority interpretations]\label{context:prop:majority}
For binary payoffs, pairwise majority of strict preferences is exactly
maximization of the approval score $v_B(o)=\sum_{c\in B}u_c(o)$.
Its lifted rule is a semilattice, and under the original diagonal
payoffs it equals $\star_B$. With general ranked payoffs, pairwise
majority can fail~\eqref{context:eq:atomic-assoc}.

For a set $B$ of $n\geq3$ distinct original owners, majority of the
proposals $(S\star_r T)_{r\in B}$ is always $S\cup T$. In fact
every quota $q<n$ gives union. Distinct tags with repeated objectives
need not have this collapse.

In a simultaneous ballot game with an odd number $n\geq3$ of distinct voters,
each voting for one of two fixed continuations, both unanimous ballot
profiles are pure Nash equilibria, whatever the voters' preferences.
Keeping all possible pure equilibrium outcomes of this ballot stage
therefore gives union of the continuation outcome sets.
\end{proposition}
\begin{proof}
For binary payoffs the signed majority margin between $x$ and $y$ is
\begin{align*}
 &\#\{c\in B:u_c(x)>u_c(y)\}-\#\{c\in B:u_c(y)>u_c(x)\}\\
 &\qquad=\sum_{c\in B}(u_c(x)-u_c(y))=v_B(x)-v_B(y).
\end{align*}
Thus the rule maximizes a scalar utility. For diagonal payoffs the
score is $[o\in B]$, recovering the shared objective.
For general rankings, let three voters have preferences
$a>b>c$, $b>c>a$, and $c>a>b$. Majority prefers $a$ to $b$, $b$ to
$c$, and $c$ to $a$. Consequently
$(\{a\}\odot\{b\})\odot\{c\}=\{c\}$, whereas
$\{a\}\odot(\{b\}\odot\{c\})=\{a\}$.

Each outcome of $S\cup T$ can be excluded by at most one original
owner's proposal. An owner excluding it must be the singleton value
of an operand. If the two operands are distinct atoms, either outcome
is still retained by its own owner; otherwise at most one owner can
absorb. Thus each outcome has at least $n-1$ approvals. In contrast,
two controller tags carrying $\star_a$ and one carrying $\star_b$
give proposal majority exactly $\star_a$.

In a unanimous ballot profile, one changed vote leaves a strict
majority unchanged, so no voter can improve. For any prescribed SPE
outcome of either continuation, prescribe it after every ballot tuple
selecting that continuation; select arbitrary SPEs after tuples
selecting the other one. Unanimous ballots for the desired branch
then give a stage Nash equilibrium. Conversely every stage outcome
comes from one continuation. This proves union for the pure outcome
observation, including when all voters strictly prefer one branch.
\end{proof}

The ballot mechanism changes the playground by inserting a simultaneous
stage, whereas the preference rule changes the node's interpretation.
The former fits the round interface of \cref{app:sealed}; its mixed
distribution semantics is richer than the pure outcome observation
just computed. For majority rules with cycles, the atomic test locates
the failure of associativity while the binary tree still specifies a
well-defined order of comparisons.

\paragraph{Disjoint teams preserve the equational theory.}
For a partition $\mathcal T=\{A_1,\ldots,A_m\}$ of $I$ into nonempty
teams, define
\[
 q(S)=\{j:S\cap A_j\ne\varnothing\},\qquad
 e(U)=\bigcup_{j\in U}A_j.
\]
The map $q$ forgets which member wins, and $e$ allows every member of
each listed team to win.

\begin{proposition}[Team refinement]\label{context:prop:team-refinement}
The maps $q$ and $e$ are homomorphisms between the winner-set algebra
with choices $\star_{A_j}$ and the ordinary $m$-role core, with
$q\circ e$ the identity. If $m\geq3$, these
algebras have exactly the same identities. This also holds when a
constant for $A_j$ is matched with the ordinary winner constant $j$.
\end{proposition}
\begin{proof}
An input lies inside $A_j$ exactly when its image under $q$ equals
$\{j\}$. If either input does, the refined output is nonempty and
lies inside $A_j$; otherwise its image is union. This proves
\[
 q(S\star_{A_j}T)=q(S)\star_j q(T).
\]
Conversely, applying~\eqref{context:eq:collective} to unions of whole
teams proves $e(U\star_j V)=e(U)\star_{A_j}e(V)$.
The constant claims and $q\circ e=\mathrm{id}$ follow from the partition.

For the other direction on identities, fix an outcome $i\in A_j$.
Map a refined value $S$ to
\[
 \pv\quad(i\in S),\qquad
 a_\ell\quad(S\subseteq A_\ell,\ \ell\ne j),\qquad
 \nv\quad\text{otherwise}.
\]
Under $\star_{A_j}$, membership of $i$ is disjunction. Under
$\star_{A_\ell}$ for $\ell\ne j$, it is disjunction unless an
operand lies inside $A_\ell$, which suppresses it. An output lies
inside an opponent team $A_\ell$ precisely when that team's choice
has an operand inside it, or both operands already lie inside it.
These are exactly the tables of the ordinary player view $C_j$
from \cref{view:thm:view}. All its states occur when $m\geq3$:
for $\nv$, use outcomes in two teams other than $A_j$.
Whole-team constants also have their usual quotient values.

Taking this homomorphism for every $i\in I$ gives an injective map
into a product of ordinary player views, since their $\pv$ coordinates
recover membership of each $i$. Thus every ordinary identity holds
in the refinement. Conversely, every identity of the refinement
holds in its quotient under $q$, the ordinary core.
\end{proof}

In particular, the nonfinite-basis theorem transfers to disjoint teams
with $m\geq3$. The three-team assumption matters. With teams
$A=\{a,b\}$ and $B=\{c\}$,
\[
 \{a\}\star_A(\{b\}\star_B\{c\})=\{a\},\qquad
 (\{a\}\star_A\{b\})\star_B(\{a\}\star_A\{c\})=\{a,b\}.
\]
So the mutual distributivity of ordinary two-player choice need not
survive refinement. Adding overlapping coalitions alongside the
individual owners is also a different extension: the proper
multi-member operations of \cref{context:prop:collective} are then
genuinely new.

\paragraph{Delegation and common admissibility.}
Another reading permits any one member of a nonempty $A$ to decide
according to its individual objective. Keeping all outcomes of all
such delegations gives
\[
 D_A(S,T)=\bigcup_{r\in A}(S\star_r T).
\]
For $A=\{r\}$ this is owned choice, but for $|A|\geq2$ it is always
$S\cup T$. Indeed, for two distinct delegates the union of their
answers is $S\cup T$: if neither choice absorbs this is immediate;
if just one absorbs, the other keeps union; if both absorb, the
inputs are their two atoms and both atoms are kept.

Taking common admissible outcomes instead gives
$\bigcap_{r\in A}(S\star_r T)$, which may be empty: with
$A=\{r,s\}$, $S=\{r\}$ and $T=\{s\}$, the two delegates insist
on different winners. This is not an ordering of the shared-objective
rule. For $A=\{a,b\}$, the inputs $\{a,b\},\{c\}$ give
common admissibility $\{a,b,c\}$ but shared-objective value
$\{a,b\}$; the inputs $\{a\},\{b\}$ give common admissibility
$\varnothing$ but the same shared-objective value $\{a,b\}$.
Thus delegation, common admissibility, and a shared objective agree
on singleton coalitions but express different collective decisions.

\subsection{Repeated descriptions with uniform component strategies}

Repeated occurrences normally have separately chosen full strategy
tables. Fixing one component SPE profile in every copy gives a different
observation in general.

\begin{proposition}[Uniform-profile observation]\label{game:prop:uniform-observation}
Let $f[x]$ be any unary owned-choice context and $G$ any closed component.
The winners of ordinary SPE profiles that use one and the same full SPE
profile of $G$ in every copy are exactly
\[
 U(f)(\den{G})=\bigcup_{g\in\den{G}}f(\{g\}).
\]
Thus uniform copies preserve the ordinary observation at $S$ precisely
when $f(S)=\bigcup_{g\in S}f(\{g\})$.
\end{proposition}
\begin{proof}
Fix the component profile. Replace every copy with the leaf bearing its
winner. Local optimality inside the component is already supplied by
that profile; local optimality at each outer node depends only on child
winners. Induction over the context therefore preserves SPE and the
selected winner. Take the union over component SPE profiles.
\end{proof}

The next theorem is the positive case for neutral union.

\begin{proposition}[Uniform copies suffice]\label{game:thm:uniform}
Fix pairwise distinct $r,s,t$, a full SPE profile $p$ of $G$,
and a full SPE profile $q$ of $H$, with
winners $g,h$. In
\[
  (G\ch r H)\ch t(G\ch s H),
\]
use exactly $p$ in both copies of $G$ and exactly $q$ in both copies of
$H$. As the three outer actions vary, the winners of ordinary SPE
profiles are exactly $\{g,h\}$. Therefore the winners of ordinary SPE
profiles satisfying this copy-equality condition are
$\den{G}\cup\den{H}$.
\end{proposition}
\begin{proof}
For any fixed three outer actions, local optimality at those nodes depends
only on the child winners $g,h$. Local optimality inside a component is
already supplied by $p,q$. Thus the expanded full profile is SPE exactly
when the corresponding profile of
$(\at g\ch r\at h)\ch t(\at g\ch s\at h)$ is SPE. By
\cref{game:thm:adequacy,game:thm:union}, its possible winners are
$\{g,h\}$. Conversely, a uniform ordinary SPE restricts to SPE profiles
in its component copies. Taking the union over $p,q$ proves the claim.
\end{proof}

We \emph{select} profiles with equal component copies from the ordinary SPE
profiles, and deviations remain unrestricted.

The term $D$ above supplies a separating case: it fixes every singleton
but adds winners on some ambiguous inputs, so $U(D)$ is the identity
while $D$ is not. Uniform observation is also not preserved by arbitrary
substitution of an expression for the input variable. The next subsection gives an explicit substitution counterexample.

\subsection{Repetition and substitution remain separate questions}

For the observation $U$ of \cref{game:prop:uniform-observation}, the
context $D$ in \eqref{game:eq:D} has $U(D)$ equal to the identity.
Nevertheless, for $H(x)=x\cup\{c\}$,
\[
 U(D\circ H)(\{a\})=I,\qquad
 (U(D)\circ U(H))(\{a\})=\{a,c\}.
\]
The uniform-profile equivalence is not preserved by substituting an
expression for its variable. Binding the entire component's selected
behavior differs from sharing only the input behavior inside several
separately solved surrounding contexts. A typed treatment of sharing
must identify that boundary explicitly.

\begin{proposition}[Unary contexts tolerating uniform selection]
For $f=\Reg(C,E)$ from the unary classification, $U(f)=f$ exactly when, for
every $k\in C$, at most one $i\ne k$ has $k\notin E_i$. There are exactly
\[
 (K+1)^K+2^K-2
\]
ordinary unary term functions preserved by $U$, whereas the number of
distinct functions in the image of $U$ is
\[
 2^{K(K-1)}+2^K-2.
\]
For $K=3$ the two counts are both $70$; for $K=4$ they are $639$ and
$4110$, respectively.
\end{proposition}
\begin{proof}
On $S$, uniform lifting returns
$S\cup\bigcup_{i\in S}E_i$. For each extra winner $k\in C$, every
ambiguous input omitting $k$ must therefore meet the rows permitting
$k$. This holds exactly when no pair of other rows both omit $k$.
For each $k$, either it is absent from $C$, or it is present and all
other rows permit it except at most one. There are $1+K$ options per
$k$. Add the nonempty constant functions and subtract the overlap $I$.

For the image count, any row family $E_i\subseteq I\setminus\{i\}$ is
realizable by setting $C=I$ in \cref{game:thm:unary}. Singleton rows
uniquely determine their union-preserving lifts. There are
$2^{K(K-1)}$ such extensive lifts. Adding constants again gives the
displayed count.
\end{proof}

For $K\geq4$, uniform lifting can even leave the ordinary unary term
functions. Choose distinct $a,b,c,d$ and a regular function with
$C=\{d\}$, $E_a=\{d\}$, and every other row empty. Its lift adds $d$
exactly when the input contains $a$. It adds $d$ at $\{a,b\}$ but not
at $\{b,c\}$, although both are ambiguous and omit $d$. No regular
ordinary term function can do that: it would need $d$ both present and
absent in its fixed extra set. This further separates the two accounts
of repetition without introducing beliefs or probabilities.

\subsection{A hidden-information obstruction}

The public recurrence does not extend unchanged to hidden information.
Suppose Alice chooses a bit and Bob guesses it, winning exactly when
he guesses correctly. With public observation, every pure SPE winner
is Bob. Hide the bit from Bob while retaining the tree and terminal
winners. At a matching pair Alice can switch her bit; at a mismatching
pair Bob can switch his guess. There is no pure Nash equilibrium.
Bob's information set crosses both branches, so neither Bob node roots
a proper continuation subgame. SPE adds no restriction beyond Nash in
this game: its refinement condition is vacuous, while Nash existence
has actually failed.

Thus winner-set equivalence for public games does not justify arbitrary
hidden composition. Sequential equilibrium introduces strategies,
beliefs, sequential rationality, and consistency, and is a different
semantic problem. 
